\chapter{Differentiation}\label{ch:b1:derivative}

\omterm{def:b1:derivative:def}{Derivatives} were computed throughout the High School volume; what was
missing is the chain of theorems that turns computation into
information about functions: Rolle's theorem, the mean value theorem,
and their consequences --- monotonicity criteria, Lipschitz bounds,
convexity. Everything in this chapter concerns functions defined on
an \omterm{prop:b1:reals:intervals}{interval} $I$.

\section{The derivative}

\begin{definition}\label{def:b1:derivative:def}
$f \colon I \to \R$ is \emph{differentiable at $x_0 \in
I$}\index{derivative} when the difference quotient
$\frac{f(x) - f(x_0)}{x - x_0}$ has a (finite) limit as $x \to x_0$;
the limit is written $f'(x_0)$. Equivalently:
\[
f(x_0 + h) = f(x_0) + f'(x_0)\,h + h\,\varepsilon(h),
\qquad \varepsilon(h) \xrightarrow[h \to 0]{} 0 ,
\]
the graph then admitting the tangent line $y = f(x_0) + f'(x_0)(x -
x_0)$. Differentiability at $x_0$ implies \omterm{def:b1:continuity:continuous}{continuity} at $x_0$ (read
the display). $f$ is differentiable on $I$ when it is at every point;
$f$ is of class $C^1$ when moreover $f'$ is \omterm{def:b1:continuity:continuous}{continuous}, and of
class $C^k$\index{Ck function@$C^k$ function} when $f$ can be
differentiated $k$ times with $f^{(k)}$ \omterm{def:b1:continuity:continuous}{continuous}.
\end{definition}

\begin{example}\label{ex:b1:derivative:notc1}
The converse of ``\omterm{def:b1:derivative:def}{differentiable} $\Rightarrow$ \omterm{def:b1:continuity:continuous}{continuous}'' fails:
$\abs{\,\cdot\,}$ at $0$. More surprisingly, \omterm{def:b1:derivative:def}{differentiable} does not
imply $C^1$: the function $f(x) = x^2 \sin\frac 1x$ ($f(0) = 0$) is
\omterm{def:b1:derivative:def}{differentiable} everywhere, with $f'(0) = 0$, but $f'(x) = 2x
\sin\frac1x - \cos\frac 1x$ has no limit at $0$
(\cref{exo:b1:derivative:2}).
\end{example}

\begin{example}[Differentiable at exactly one point]\label{ex:b1:derivative:onepoint}
Let $f(x) = x^2$ for $x \in \Q$ and $f(x) = 0$ for $x \notin
\Q$. At $0$: $\bigl|\frac{f(h) - 0}{h}\bigr| \leq \abs h \to 0$,
so $f$ is \omterm{def:b1:derivative:def}{differentiable} at $0$ with $f'(0) = 0$. At any $x_0
\neq 0$, $f$ is not even \omterm{def:b1:continuity:continuous}{continuous}: rational and irrational
sequences converging to $x_0$ send $f$ to $x_0^2 \neq 0$ and to
$0$ respectively (\omterm{def:b1:topology:dense}{density}, \cref{thm:b1:reals:density}). So
differentiability is a genuinely \emph{pointwise} notion: it can
hold at one point of $\R$ and nowhere else. The moral for
practice: \omterm{def:b1:logic:statement}{statements} like the monotonicity criterion or Rolle
require the \omterm{def:b1:derivative:def}{derivative} \emph{on an \omterm{prop:b1:reals:intervals}{interval}} --- possessing
$f'(x_0)$ at isolated points, however many, supports no global
conclusion whatsoever.
\end{example}

\begin{theorem}[Operations]\label{thm:b1:derivative:operations}
If $f, g$ are \omterm{def:b1:derivative:def}{differentiable} at $x_0$ (and where the formulas make
sense):
\[
(f + g)' = f' + g', \qquad
(fg)' = f'g + fg', \qquad
\Bigl(\frac fg\Bigr)' = \frac{f'g - fg'}{g^2},
\]
and if $g$ is \omterm{def:b1:derivative:def}{differentiable} at $f(x_0)$:
$\;(g \circ f)'(x_0) = g'\bigl(f(x_0)\bigr)\, f'(x_0)$ (chain
rule\index{chain rule}).
\end{theorem}

\begin{proof}
Sum: immediate. Product: write
\[
f(x)g(x) - f(x_0)g(x_0)
= \bigl(f(x) - f(x_0)\bigr) g(x) + f(x_0)\bigl(g(x) - g(x_0)\bigr),
\]
divide by $x - x_0$ and let $x \to x_0$ ($g$ is \omterm{def:b1:continuity:continuous}{continuous} at $x_0$).
Quotient: treat $\frac 1g$ via $\frac{1/g(x) - 1/g(x_0)}{x - x_0} =
\frac{-1}{g(x)g(x_0)}\cdot\frac{g(x) - g(x_0)}{x - x_0}$, then apply
the product rule. Chain rule: with $y_0 = f(x_0)$, define
$\theta(y) = \frac{g(y) - g(y_0)}{y - y_0}$ for $y \neq y_0$ and
$\theta(y_0) = g'(y_0)$: $\theta$ is \omterm{def:b1:continuity:continuous}{continuous} at $y_0$, and for
$x \neq x_0$,
\[
\frac{g(f(x)) - g(f(x_0))}{x - x_0}
= \theta\bigl(f(x)\bigr)\cdot \frac{f(x) - f(x_0)}{x - x_0}
\longrightarrow g'(y_0)\, f'(x_0),
\]
the first factor by composition of limits (this device handles the
case $f(x) = f(x_0)$ cleanly, where the naive ``multiply and divide
by $f(x) - f(x_0)$'' breaks).
\end{proof}

\begin{theorem}[Derivative of an inverse function]\label{thm:b1:derivative:inverse}
Let $f$ be \omterm{def:b1:continuity:continuous}{continuous} and strictly monotonic on $I$, \omterm{def:b1:derivative:def}{differentiable}
at $x_0$ with $f'(x_0) \neq 0$. Then $f^{-1}$
(\cref{thm:b1:continuity:bijection}) is \omterm{def:b1:derivative:def}{differentiable} at $y_0 =
f(x_0)$, with
\[
(f^{-1})'(y_0) = \frac{1}{f'(x_0)} = \frac{1}{f'\bigl(f^{-1}(y_0)\bigr)} .
\]
If $f'(x_0) = 0$, the inverse has a vertical tangent at $y_0$.
\end{theorem}

\begin{proof}
For $y \to y_0$, set $x = f^{-1}(y)$: \omterm{def:b1:continuity:continuous}{continuity} of $f^{-1}$ gives $x
\to x_0$, and
\[
\frac{f^{-1}(y) - f^{-1}(y_0)}{y - y_0}
= \frac{x - x_0}{f(x) - f(x_0)}
= \frac{1}{\dfrac{f(x) - f(x_0)}{x - x_0}}
\longrightarrow \frac{1}{f'(x_0)} .
\]
Vertical-tangent claim: if $f'(x_0) = 0$, the displayed quotient
is the reciprocal of a quantity that tends to $0$ while keeping
one constant sign (for $f$ strictly increasing, $\frac{f(x) -
f(x_0)}{x - x_0} > 0$ for all $x \neq x_0$): the difference
quotient of $f^{-1}$ therefore tends to $+\infty$ (to $-\infty$
for $f$ decreasing). The inverse remains \omterm{def:b1:continuity:continuous}{continuous} but is not
\omterm{def:b1:derivative:def}{differentiable} at $y_0$ --- its graph, the reflection of $f$'s
across the diagonal, stands vertical exactly where $f$'s ran
horizontal, as $x^{1/3}$ at $0$ illustrates against $x^3$.
\end{proof}

\begin{example}[Inverse derivatives, twice]\label{ex:b1:derivative:inverseaction}
The theorem recomputes the classical \omterm{def:b1:derivative:def}{derivatives} with no limit
work. For $\ln = \exp^{-1}$: at $y = \eu^x$,
\[
(\ln)'(y) = \frac{1}{\exp'(x)} = \frac{1}{\eu^{x}} = \frac1y ,
\]
valid for every $y > 0$ since $\exp' = \exp$ never vanishes. For
$\arctan = \tan^{-1}$: at $y = \tan x$,
\[
(\arctan)'(y) = \frac{1}{1 + \tan^2 x} = \frac{1}{1 + y^2} ,
\]
using $\tan' = 1 + \tan^2 > 0$. The closing insight: the formula
converts knowledge about a function into knowledge about its
inverse at the price of one substitution --- and the
substitution ($x = \ln y$, $x = \arctan y$) is exactly the
\omterm{def:b1:logic:statement}{statement} that the two variables live on opposite sides of the
bijection.
\end{example}

\section{Rolle and the mean value theorem}

\begin{proposition}[Interior extremum]\label{prop:b1:derivative:fermat}
If $f$ is \omterm{def:b1:derivative:def}{differentiable} at an \emph{\omterm{def:b1:topology:closure}{interior}} point $x_0$ of $I$ and
has a local extremum there, then $f'(x_0) = 0$.
\end{proposition}

\begin{proof}
Say a local maximum: there is $r > 0$ with $f(x) \leq f(x_0)$
for $\abs{x - x_0} \leq r$, and interiority guarantees both
sides of $x_0$ are available within $I$. For $0 < h \leq r$ the
quotient $\frac{f(x_0 + h) - f(x_0)}{h}$ has numerator $\leq 0$
and denominator $> 0$: it is $\leq 0$, and its limit $f'(x_0)$
inherits $\leq 0$ (wide inequalities pass to limits,
\cref{thm:b1:seq:order}); for $-r \leq h < 0$ the quotient is
$\geq 0$, giving $f'(x_0) \geq 0$. Hence $f'(x_0) = 0$. (At an
endpoint, only one sign is available: the conclusion fails there ---
think of $x$ on $\intcc{0}{1}$, maximal at $1$ with \omterm{def:b1:derivative:def}{derivative}
$1$.)
\end{proof}

\begin{theorem}[Rolle]\label{thm:b1:derivative:rolle}
Let $f$ be \omterm{def:b1:continuity:continuous}{continuous} on $\intcc{a}{b}$, \omterm{def:b1:derivative:def}{differentiable} on
$\intoo{a}{b}$, with $f(a) = f(b)$. Then $f'(c) = 0$ for some $c \in
\intoo{a}{b}$.\index{Rolle's theorem}
\end{theorem}

\begin{proof}
By the extreme value theorem (\cref{thm:b1:continuity:evt}), $f$
attains its maximum and minimum on $\intcc{a}{b}$. If both are
attained at endpoints, then (since $f(a) = f(b)$) max $=$ min and $f$
is constant: any \omterm{def:b1:topology:closure}{interior} $c$ works. Otherwise an extremum is
attained at an \omterm{def:b1:topology:closure}{interior} point $c$, and
\cref{prop:b1:derivative:fermat} gives $f'(c) = 0$.
\end{proof}

\begin{theorem}[Mean value theorem]\label{thm:b1:derivative:mvt}
Let $f$ be \omterm{def:b1:continuity:continuous}{continuous} on $\intcc{a}{b}$, \omterm{def:b1:derivative:def}{differentiable} on
$\intoo{a}{b}$. There exists $c \in \intoo{a}{b}$ with
\[
f(b) - f(a) = f'(c)\,(b - a) .
\]
\emph{Mean value inequality:} if moreover $m \leq f' \leq M$ on
$\intoo{a}{b}$, then $m(b-a) \leq f(b) - f(a) \leq M(b-a)$; in
particular $\abs{f'} \leq K$ implies that $f$ is $K$-Lipschitz.
\index{mean value theorem}
\end{theorem}

\begin{proof}
Apply Rolle to $g(x) = f(x) - \frac{f(b) - f(a)}{b - a}(x - a)$: $g$
is \omterm{def:b1:continuity:continuous}{continuous} on $\intcc{a}{b}$, \omterm{def:b1:derivative:def}{differentiable} inside, and $g(a) =
f(a) = g(b)$. At the point $c$ where $g'(c) = 0$: $f'(c) =
\frac{f(b)-f(a)}{b-a}$. The inequality follows by bounding $f'(c)$;
the Lipschitz \omterm{def:b1:logic:statement}{statement} applies it to every pair of points.
\end{proof}

\begin{omfigure}
\begin{tikzpicture}
\begin{axis}[omaxis, width=8.2cm, height=5.4cm,
    xmin=0, xmax=4.6, ymin=0, ymax=3.2,
    xtick={0.7,2.2,3.9}, xticklabels={$a$,$c$,$b$}, ytick=\empty]
  \addplot[omDef, very thick, domain=0.7:3.9, samples=100]
    {0.5 + 1.4*ln(x+0.3)};
  \addplot[omExo, thick] coordinates {(0.7,1.5) (3.9,2.5)};
  \addplot[omThm, thick, dashed, domain=1.1:3.4]
    {2.06 + 0.3125*(x - 2.2)};
  \fill[omThm] (axis cs:2.2,2.06) circle (1.6pt);
\end{axis}
\end{tikzpicture}

{\small The mean value theorem: some tangent (dashed) is parallel to
the chord (gray). Its abscissa $c$ is where Rolle's theorem, applied
to the function minus its chord, finds a critical point.}
\end{omfigure}

\begin{example}[Newton's method is Heron's]\label{ex:b1:derivative:newton}
Newton's method for solving $f(x) = 0$ replaces the curve by
its tangent at the current guess $x_n$ and takes the tangent's
root as the next guess:
\[
0 = f(x_n) + f'(x_n)(x_{n+1} - x_n)
\quad\Longrightarrow\quad
x_{n+1} = x_n - \frac{f(x_n)}{f'(x_n)} .
\]
Run it on $f(x) = x^2 - 2$:
\[
x_{n+1} = x_n - \frac{x_n^2 - 2}{2x_n}
= \frac{x_n}{2} + \frac{1}{x_n}
= \frac12\Bigl(x_n + \frac{2}{x_n}\Bigr) :
\]
exactly Heron's iteration (\cref{ex:b1:seq:heron}), two
millennia early. The quadratic speed observed there is now
explained by the tangent picture: near a simple root, curve and
tangent differ by a second-order error, so each step roughly
squares the error --- the general \omterm{def:b1:logic:statement}{statement} follows from the
Taylor bounds of \cref{ch:b1:taylor}. The closing insight:
where dichotomy (\cref{ex:b1:continuity:dichotomy}) uses only
\omterm{def:b1:continuity:continuous}{continuity} and gains one bit per step, Newton spends a
\omterm{def:b1:derivative:def}{derivative} to double the number of correct digits per step.
\end{example}

\begin{example}[The mean value theorem as an estimator]\label{ex:b1:derivative:mvtestimate}
How large is $\sqrt{101}$? Apply the theorem to $f(t) = \sqrt t$
on $\intcc{100}{101}$: for some $c \in \intoo{100}{101}$,
\[
\sqrt{101} - 10 = \frac{1}{2\sqrt c},
\qquad\text{so}\qquad
\frac{1}{2\sqrt{101}} < \sqrt{101} - 10 < \frac{1}{20} = 0.05 ,
\]
and since $\sqrt{101} < 10.05$, the left bound exceeds
$\frac{1}{20.1} > 0.0497$: thus $10.0497 < \sqrt{101} < 10.05$
(true value $10.049875\dots$) --- three correct decimals from one
\omterm{def:b1:derivative:def}{derivative} evaluation. Likewise $\abs{\sin a - \sin b} \leq
\abs{a - b}$ (bound $\abs{\cos}\leq 1$): the Lipschitz estimates
used since \cref{ch:b1:seq} are all this theorem. The closing
insight: the mean value theorem is a zeroth-order Taylor formula
--- it trades one unknown point $c$ for a hard inequality, and
\cref{ch:b1:taylor} will iterate exactly this trade.
\end{example}

\begin{corollary}[Monotonicity criterion]\label{cor:b1:derivative:monotone}
Let $f$ be \omterm{def:b1:continuity:continuous}{continuous} on $I$, \omterm{def:b1:derivative:def}{differentiable} on the \omterm{def:b1:topology:closure}{interior}.
\begin{enumerate}
  \item $f' \geq 0$ on the \omterm{def:b1:topology:closure}{interior} $\iff$ $f$ is increasing; $f' = 0$
        $\iff$ $f$ constant.
  \item If $f' > 0$ except at finitely many points where it vanishes,
        $f$ is \emph{strictly} increasing.
\end{enumerate}
\end{corollary}

\begin{proof}
If $f' \geq 0$: for $x < y$ in $I$, the mean value theorem on
$\intcc{x}{y}$ gives $f(y) - f(x) = f'(c)(y - x) \geq 0$. Conversely,
difference quotients of an increasing function are $\geq 0$, so their
limits are too. The constant case: apply the previous to $f'$ and
$-f' \geq 0$. Strict version: $f$ is increasing; equality $f(x) =
f(y)$ for $x < y$ would freeze $f$ on $\intcc{x}{y}$, forcing $f' =
0$ there --- infinitely many points.
\end{proof}

\begin{example}[Equal derivatives, unequal functions]\label{ex:b1:derivative:intervalmatters}
On $\R^* = \intoo{-\infty}{0} \cup \intoo{0}{+\infty}$, both
$f(x) = \ln\abs x$ and $g(x) = \ln\abs x + \mathbf{1}_{x>0}$
(add $1$ on the right half-line only) satisfy $f' = g' =
\frac1x$. They do not differ by a constant: the criterion
``$f' = 0 \implies f$ constant'' is an \emph{\omterm{prop:b1:reals:intervals}{interval}} \omterm{def:b1:logic:statement}{statement}
--- its proof runs the mean value theorem between two points,
which requires the whole segment joining them to lie in the
domain. On each half-line separately, the primitives of
$\frac1x$ are $\ln\abs x + c$ with one constant per half-line,
two independent constants in total. \Cref{ch:b1:integration}
inherits this fine print: ``the'' primitive of a function is
well defined up to a constant \emph{on each \omterm{prop:b1:reals:intervals}{interval} of its
domain}, and antiderivative tables silently assume connectivity.
\end{example}

\begin{example}[Strictness for free]\label{ex:b1:derivative:strictness}
$x \mapsto x^3$ is \emph{strictly} increasing on $\R$ even
though its \omterm{def:b1:derivative:def}{derivative} vanishes at $0$: the criterion's clause
``$f' > 0$ except at finitely many points'' is exactly designed
for such flat points. By contrast, $f' \geq 0$ alone only gives
increase in the wide sense (a constant function qualifies), and
a \omterm{def:b1:derivative:def}{derivative} vanishing on a whole subinterval does freeze the
function there. The practical rule: to claim strict
monotonicity, list the zeros of $f'$; finitely many (or more
generally, none on any subinterval) is harmless, an \omterm{prop:b1:reals:intervals}{interval} of
them is fatal.
\end{example}

\begin{example}[A complete variation study]\label{ex:b1:derivative:variations}
Study $f(x) = x^3 - 3x + 1$ on $\R$. \omterm{def:b1:derivative:def}{Derivative}: $f'(x) =
3(x^2 - 1)$, positive on $\intoo{-\infty}{-1}$, negative on
$\intoo{-1}{1}$, positive on $\intoo{1}{+\infty}$: by the
monotonicity criterion, $f$ increases, then decreases, then
increases, with a local maximum $f(-1) = 3$ and a local minimum
$f(1) = -1$. Limits: $\mp\infty$ at $\mp\infty$. Consequences,
read off the variation table with the intermediate value
theorem on each monotone branch: $f$ vanishes exactly once in
each of
\[
\intoo{-\infty}{-1}, \qquad \intoo{-1}{1}, \qquad
\intoo{1}{+\infty}
\]
(the values at the junctions have opposite signs: $3 > 0 > -1$),
so the equation $x^3 - 3x + 1 = 0$ has exactly three real
roots; numerically they sit near $-1.88$, $0.35$, $1.53$. The
closing insight: a variation table is a \emph{proof device},
not a sketch --- monotone branch plus sign change equals
exactly one root, and the table enumerates the branches
exhaustively.
\end{example}

\begin{theorem}[Leibniz formula]\label{thm:b1:derivative:leibniz}
If $f, g$ are $n$ times \omterm{def:b1:derivative:def}{differentiable}, so is $fg$,
and\index{Leibniz formula}
\[
(fg)^{(n)} = \sum_{k=0}^{n} \binom nk f^{(k)}\, g^{(n-k)} .
\]
\end{theorem}

\begin{proof}
Induction on $n$, exactly parallel to the binomial theorem. The
case $n = 1$ is the product rule. Assuming the formula at rank
$n$, differentiate once more:
\[
(fg)^{(n+1)} = \sum_{k=0}^{n} \binom nk
\Bigl( f^{(k+1)} g^{(n-k)} + f^{(k)} g^{(n-k+1)} \Bigr),
\]
then reindex the first sum with $j = k + 1$ and collect the
coefficient of $f^{(j)} g^{(n+1-j)}$: it is $\binom{n}{j-1} +
\binom nj = \binom{n+1}{j}$ by Pascal's rule
(\cref{prop:b1:counting:identities}), the \omterm{def:b1:topology:closure}{boundary} terms $j = 0$
and $j = n + 1$ carrying $\binom{n+1}{0} = \binom{n+1}{n+1} =
1$ as they should.
\end{proof}

\begin{example}[Leibniz in action]\label{ex:b1:derivative:leibnizaction}
Compute $\bigl(x^2 \eu^x\bigr)^{(n)}$ for $n \geq 2$. Take $f =
x^2$, whose \omterm{def:b1:derivative:def}{derivatives} die quickly ($f' = 2x$, $f'' = 2$,
$f^{(k)} = 0$ for $k \geq 3$), and $g = \eu^x$: only three terms
of the Leibniz sum survive,
\[
\bigl(x^2\eu^x\bigr)^{(n)}
= \binom n0 x^2 \eu^x + \binom n1 (2x)\,\eu^x + \binom n2\,
2\,\eu^x
= \eu^x\bigl(x^2 + 2nx + n(n-1)\bigr).
\]
Sanity check at $n = 1$: $\eu^x(x^2 + 2x)$, which is indeed
$(x^2\eu^x)'$. The closing insight: use Leibniz when one factor
is a \omterm{def:b1:poly:def}{polynomial} --- the sum then has only $\deg + 1$ terms, and
the formula is a closed form, not an abstract identity. (For two
infinitely-lively factors like $\eu^x\sin x$, complex
exponentials from \cref{ch:b1:complex} are the better tool.)
\end{example}

\section{Convexity}

\begin{definition}\label{def:b1:derivative:convex}
$f \colon I \to \R$ is \emph{convex}\index{convex function} when
every chord lies above the graph:
\[
\forall x, y \in I,\ \forall t \in \intcc{0}{1}, \quad
f\bigl(tx + (1-t)y\bigr) \leq t f(x) + (1-t) f(y).
\]
($f$ is \emph{concave} when $-f$ is convex.)
\end{definition}

\begin{theorem}[Differential characterizations]\label{thm:b1:derivative:convexchar}
Let $f$ be \omterm{def:b1:derivative:def}{differentiable} on $I$. The following are equivalent:
\begin{enumerate}
  \item $f$ is \omterm{def:b1:derivative:convex}{convex};
  \item $f'$ is increasing on $I$;
  \item the graph lies above every tangent: $f(y) \geq f(x) +
        f'(x)(y - x)$ for all $x, y \in I$.
\end{enumerate}
If $f$ is twice \omterm{def:b1:derivative:def}{differentiable}: $f$ \omterm{def:b1:derivative:convex}{convex} $\iff f'' \geq 0$.
\end{theorem}

\begin{proof}
(1 $\Rightarrow$ 3) Convexity written as $\frac{f(x + t(y-x)) -
f(x)}{t} \leq f(y) - f(x)$ for $t \in \intoc{0}{1}$; let $t \to 0^+$:
$f'(x)(y - x) \leq f(y) - f(x)$.

(3 $\Rightarrow$ 2) For $x < y$, the two tangent inequalities at $x$
and at $y$ give $f'(x)(y-x) \leq f(y) - f(x) \leq f'(y)(y - x)$,
hence $f'(x) \leq f'(y)$.

(2 $\Rightarrow$ 1) Fix $x < y$ and $t \in \intoo{0}{1}$, and let $z
= tx + (1-t)y \in \intoo{x}{y}$. By the mean value theorem on
$\intcc{x}{z}$ and $\intcc{z}{y}$: there are $c_1 < z < c_2$ with
\[
\frac{f(z) - f(x)}{z - x} = f'(c_1) \leq f'(c_2)
= \frac{f(y) - f(z)}{y - z} ,
\]
and clearing denominators ($z - x = (1-t)(y-x)$, $y - z = t(y-x)$)
rearranges exactly into the convexity inequality.

Twice \omterm{def:b1:derivative:def}{differentiable} case: $f'' \geq 0 \iff f'$ increasing
(\cref{cor:b1:derivative:monotone}).
\end{proof}

\begin{omfigure}
\begin{tikzpicture}
\begin{axis}[omaxis, width=8.2cm, height=5.6cm,
    xmin=-2.4, xmax=2.6, ymin=-0.8, ymax=3.4,
    xtick=\empty, ytick=\empty]
  \addplot[omDef, very thick, domain=-2.1:2.3, samples=100]
    {0.5*x*x};
  \addplot[omExo, thick] coordinates {(-1.6,1.28) (2,2)};
  \addplot[omThm, thick, dashed, domain=-1.2:2.4] {0.72*(x-0.72)+0.2592};
\end{axis}
\end{tikzpicture}

{\small Convexity, twice: every chord (gray) lies above the graph,
and the graph lies above every tangent (dashed).}
\end{omfigure}

\begin{example}[Classical convexity inequalities]\label{ex:b1:derivative:convexineq}
$\exp$ is \omterm{def:b1:derivative:convex}{convex} ($\exp'' = \exp > 0$): its tangent at $0$ gives
$\eu^x \geq 1 + x$ for all $x$. $\ln$ is concave: its tangent at $1$
gives $\ln x \leq x - 1$; its chords give, for $0 < a \leq b$, the
inequality between geometric and arithmetic means: taking $t =
\frac12$ in concavity,
\[
\ln\frac{a + b}{2} \geq \frac{\ln a + \ln b}{2} = \ln\sqrt{ab},
\qquad\text{so}\qquad
\sqrt{ab} \leq \frac{a+b}{2} .
\]
The general arithmetic--geometric inequality is
\cref{exo:b1:derivative:9}.
\end{example}

\begin{example}[A convexity inequality from scratch]\label{ex:b1:derivative:xlnx}
The function $f(t) = t\ln t$ is \omterm{def:b1:derivative:convex}{convex} on $\intoo{0}{+\infty}$:
$f''(t) = \frac1t > 0$. Its midpoint inequality, multiplied by
$2$, reads: for all $a, b > 0$,
\[
a\ln a + b\ln b \;\geq\; (a + b)\,\ln\frac{a + b}{2} ,
\]
with equality iff $a = b$ (strict convexity). Test drive: $a =
1$, $b = 3$ gives $3\ln 3 = 3.296$ against $4\ln 2 = 2.773$.
This innocuous inequality is the two-point case of the
\emph{entropy} comparison that reappears with Jensen's
inequality (\cref{exo:b1:derivative:9}) and in the
information-theoretic asymptotics of the Year 3 volume. The
closing insight: to manufacture an inequality, find a function
whose second \omterm{def:b1:derivative:def}{derivative} has a sign and write down what convexity
says --- the differential characterization turns one sign check
into infinitely many inequalities.
\end{example}

\begin{remark}[Common pitfalls with derivatives]
(i) \emph{A positive \omterm{def:b1:derivative:def}{derivative} at one point does not give
monotonicity near it}: $f(x) = \frac x2 + x^2\sin\frac1x$ (with
$f(0) = 0$) has $f'(0) = \frac12 > 0$, yet
\[
f'(x) = \frac12 + 2x\sin\frac1x - \cos\frac1x
\]
equals $-\frac12$ at each $x_n = \frac{1}{2\pi n}$: every
\omterm{def:b1:topology:open}{neighborhood} of $0$ contains descents. Monotonicity needs $f'
\geq 0$ \emph{on an \omterm{prop:b1:reals:intervals}{interval}}
(\cref{cor:b1:derivative:monotone}); the pointwise sign only
controls the crossing of the tangent line. (ii) \emph{Rolle's
three hypotheses are all active}: $\abs x$ on $\intcc{-1}{1}$
(no \omterm{def:b1:topology:closure}{interior} differentiability), $x$ on $\intcc{0}{1}$ (ends not
equal), and $x - \lfloor x\rfloor$ on $\intcc{0}{1}$ (\omterm{def:b1:continuity:continuous}{continuity}
fails at $1$) each break exactly one hypothesis and the
conclusion. (iii) \emph{\omterm{def:b1:derivative:def}{Derivatives} may be discontinuous, but
not arbitrarily}: $f'$ can oscillate
(\cref{ex:b1:derivative:notc1}) yet always satisfies the
intermediate value property (Darboux,
\cref{exo:b1:derivative:10}): a \omterm{def:b1:derivative:def}{derivative} never jumps --- if
you compute a one-sided ``\omterm{def:b1:derivative:def}{derivative} limit'' with a jump, you
have differentiated a \omterm{def:b1:derivative:def}{non-differentiable} function. (iv) \emph{The
inverse formula needs $f' \neq 0$}: $x \mapsto x^3$ is a smooth
strictly increasing bijection whose inverse $x^{1/3}$ has a
vertical tangent at $0$ --- differentiability of the inverse is
lost exactly where $f'$ vanishes
(\cref{thm:b1:derivative:inverse}).
\end{remark}

\begin{remark}[Where the mean value theorem works next]
Nearly every quantitative \omterm{def:b1:logic:statement}{statement} of the next chapters is this
chapter's mean value theorem in costume: the fundamental theorem
of calculus (\cref{ch:b1:integration}) differentiates the area
function and concludes with the monotonicity criterion; the
Taylor--Lagrange formula (\cref{ch:b1:taylor}) is the mean value
theorem iterated $n$ times; the error analysis of Newton's method
and of fixed-point iterations (\cref{exo:b1:derivative:11}) is
the Lipschitz form; and this chapter's weekend problem
(\cref{pb:b1:derivative:1}) turns the same Lipschitz bound into
number theory --- a repulsion inequality between \omterm{pb:b1:logic:1}{algebraic
numbers} and rationals, yielding the first \omterm{pb:b1:logic:1}{transcendental number}
in history. In the Year 2 volume, the mean value
\emph{inequality} survives in several variables when the equality
does not.
\end{remark}

\begin{example}[Young's inequality from concavity]\label{ex:b1:derivative:young}
Let $p, q > 1$ with $\frac1p + \frac1q = 1$. For all $a, b > 0$:
\[
ab \;\leq\; \frac{a^p}{p} + \frac{b^q}{q} .
\]
Proof by one application of the concavity of $\ln$ with weights
$\frac1p, \frac1q$ (the two-point Jensen inequality, as in
\cref{exo:b1:derivative:9}):
\[
\ln\Bigl(\frac{a^p}{p} + \frac{b^q}{q}\Bigr)
\;\geq\; \frac1p \ln(a^p) + \frac1q \ln(b^q)
= \ln a + \ln b = \ln(ab),
\]
and $\ln$ increasing converts the inequality of logarithms into
the claim; equality iff $a^p = b^q$ (strict concavity). The case
$p = q = 2$ is the arithmetic-geometric inequality $ab \leq
\frac{a^2 + b^2}{2}$ in disguise. The closing insight: Young's
inequality is the algebraic seed of the H\"older and Minkowski
inequalities of the Year 2 volume --- one concavity \omterm{def:b1:logic:statement}{statement}
about $\ln$, harvested for norms.
\end{example}

\begin{remark}[Perspectives inside this volume]
The \omterm{def:b1:derivative:def}{derivative} acquires three new lives before the volume ends.
In \cref{ch:b1:taylor} it iterates: $n$ \omterm{def:b1:derivative:def}{derivatives} at a point
compress into one \omterm{def:b1:poly:def}{polynomial} plus a controlled error, and the
mean value theorem becomes the Lagrange remainder. In
\cref{ch:b1:curves}, differentiation turns geometric: for a
parametrized curve $t \mapsto (x(t), y(t))$, the pair $(x'(t),
y'(t))$ is a velocity \emph{vector}, tangency becomes
collinearity, and critical points become cusps to classify. In
\cref{ch:b1:multivar}, one variable is frozen at a time: partial
\omterm{def:b1:derivative:def}{derivatives} repeat this chapter twice over, and the tangent line
grows into a tangent plane. All three chapters inherit the same
grammar --- local linear approximation plus an error term ---
first spoken here.
\end{remark}

\section{Exercises}

\begin{exercise}[$\star$]\label{exo:b1:derivative:1}
Differentiate (specifying the domains): $x^x$;
$\;\ln\bigl(x + \sqrt{x^2+1}\bigr)$; $\;\arctan\frac{1}{x}$;
$\;\sqrt{1 + \eu^{2x}}$.
\end{exercise}

\begin{exercise}[$\star$]\label{exo:b1:derivative:2}
Complete \cref{ex:b1:derivative:notc1}: prove that $f(x) = x^2
\sin\frac1x$, $f(0) = 0$, is \omterm{def:b1:derivative:def}{differentiable} at $0$ with $f'(0) = 0$,
and that $f'$ has no limit at $0$.
\end{exercise}

\begin{exercise}[$\star$]\label{exo:b1:derivative:3}
Using the mean value theorem or the tangent inequalities, prove that
for all $x > 0$:
\[
\frac{x}{1 + x} < \ln(1 + x) < x .
\]
Deduce $\lim_{n\to\infty} \bigl(1 + \frac xn\bigr)^n = \eu^x$ for
every $x > 0$.
\end{exercise}

\begin{exercise}[$\star$]\label{exo:b1:derivative:4}
Let $P$ be a real \omterm{def:b1:poly:def}{polynomial} with $k$ distinct real roots. Prove that
$P'$ has at least $k - 1$ distinct real roots, interlaced with those
of $P$. Deduce that if $P$ has all its roots real, so does $P'$.
\end{exercise}

\begin{exercise}[$\star\star$]\label{exo:b1:derivative:5}
Let $f$ be \omterm{def:b1:derivative:def}{differentiable} on $\R$ with $f' (x)\to \ell$ as $x \to
+\infty$. Prove that $\frac{f(x)}{x} \to \ell$ \emph{(mean value
theorem on $\intcc{A}{x}$)}. Does $f(x+1) - f(x) \to \ell$ hold too?
\end{exercise}

\begin{exercise}[$\star\star$]\label{exo:b1:derivative:6}
(A discrete Rolle) Let $f$ be $n$ times \omterm{def:b1:derivative:def}{differentiable} on $I$ and
vanish at $n + 1$ distinct points. Prove that $f^{(n)}$ vanishes at
least once. Application: a \omterm{def:b1:poly:def}{polynomial} of degree $\leq n$ vanishing at
$n+1$ points is zero (again).
\end{exercise}

\begin{exercise}[$\star\star$]\label{exo:b1:derivative:7}
Let $f$ be twice \omterm{def:b1:derivative:def}{differentiable} on $\intcc{a}{b}$ with $f(a) = f(b)
= 0$ and $f(x_0) > 0$ for some \omterm{def:b1:topology:closure}{interior} $x_0$. Prove that $f''(c) <
0$ for some $c \in \intoo{a}{b}$. \emph{(Two mean value theorems and
a comparison of slopes.)}
\end{exercise}

\begin{exercise}[$\star\star$]\label{exo:b1:derivative:8}
Study the function $f(x) = \dfrac{\ln x}{x}$ on
$\intoo{0}{+\infty}$: variations, limits, maximum. Deduce that $a^b
> b^a$ for all reals $\eu \leq a < b$, and settle the famous special
case: which of $\eu^\pi$, $\pi^\eu$ is larger? Check against the
small integer pairs $(2,3)$ and $(2,4)$: why do they behave
differently?
\end{exercise}

\begin{exercise}[$\star\star$]\label{exo:b1:derivative:9}
(Arithmetic--geometric inequality) Using the concavity of $\ln$ with
general weights (Jensen's inequality for $n$ points, to be proved by
induction on $n$), show that for positive reals $a_1, \dots, a_n$:
\[
\sqrt[n]{a_1 a_2 \cdots a_n} \leq \frac{a_1 + \dots + a_n}{n},
\]
with equality iff all $a_i$ are equal.
\end{exercise}

\begin{exercise}[$\star\star\star$]\label{exo:b1:derivative:10}
(Darboux: \omterm{def:b1:derivative:def}{derivatives} take intermediate values) Let $f$ be
\omterm{def:b1:derivative:def}{differentiable} on $I$ and $a < b$ in $I$ with $f'(a) < v < f'(b)$.
By considering $g(x) = f(x) - vx$ and the point where $g$ attains
its minimum on $\intcc{a}{b}$, prove that $f'(c) = v$ for some $c
\in \intoo{a}{b}$ --- even though $f'$ need not be \omterm{def:b1:continuity:continuous}{continuous}
(\cref{exo:b1:derivative:2}).
\end{exercise}

\begin{exercise}[$\star\star\star$]\label{exo:b1:derivative:11}
Let $f \colon \R \to \R$ be \omterm{def:b1:derivative:def}{differentiable} with $\abs{f'(x)} \leq k
< 1$ for all $x$ (a \emph{contraction}). Prove that $f$ has exactly
one fixed point $\ell$, and that every sequence $u_{n+1} = f(u_n)$
converges to $\ell$ with $\abs{u_n - \ell} \leq k^n\abs{u_0 -
\ell}$. \emph{(Existence: apply the intermediate value theorem to
$f(x) - x$ on a large segment, using the Lipschitz bound; or use
completeness with the Cauchy criterion.)}
\end{exercise}

\begin{exercise}[$\star\star\star$]\label{exo:b1:derivative:12}
(Cauchy's mean value theorem and l'Hospital's rule)
\begin{enumerate}
  \item Let $f, g$ be \omterm{def:b1:continuity:continuous}{continuous} on $\intcc{a}{b}$,
        \omterm{def:b1:derivative:def}{differentiable} on $\intoo{a}{b}$, with $g'$ never zero
        there. Prove that $g(b) \neq g(a)$ and that some $c \in
        \intoo{a}{b}$ satisfies
        \[
        \frac{f(b) - f(a)}{g(b) - g(a)} = \frac{f'(c)}{g'(c)}
        \]
        \emph{(apply Rolle to $h = f - \lambda g$ for the right
        constant $\lambda$)}.
  \item Deduce l'Hospital's rule in the $\frac00$ form at a
        point: if $f(a) = g(a) = 0$ and
        $\frac{f'(x)}{g'(x)} \to \ell$ as $x \to a^+$, then
        $\frac{f(x)}{g(x)} \to \ell$.
  \item Show the converse fails: for $f(x) = x^2\sin\frac1x$
        ($f(0) = 0$) and $g(x) = x$, the quotient
        $\frac{f}{g}$ has a limit at $0$ but
        $\frac{f'}{g'}$ has none.
\end{enumerate}
\end{exercise}

\section{Problem: Liouville's inequality and the first
transcendental number}

\begin{problem}[{Weekend problem --- \omterm{pb:b1:logic:1}{algebraic numbers} repel
rationals: $\abs{x - p/q} \geq C/q^d$, and the transcendence of
$\sum 10^{-n!}$}]
\label{pb:b1:derivative:1}
A real number is \emph{algebraic} when it is a root of a nonzero
\omterm{def:b1:poly:def}{polynomial} with integer coefficients, and \emph{transcendental}
otherwise. In 1844 Liouville produced the first number ever
\emph{proved} transcendental, and the engine of his proof is this
chapter's mean value theorem: an \omterm{pb:b1:logic:1}{algebraic number} of degree $d$
cannot be approximated by rationals better than $C/q^d$ --- so a
number approximable \emph{faster than every power} cannot be
algebraic. This problem builds the inequality, constructs
Liouville's number $L = 0.110001000\dots$ (ones at the factorial
positions, via the digit machinery of \cref{pb:b1:reals:1}),
proves its transcendence, and ends with Cantor's rival proof and
effective bounds for $\sqrt2$ and
$2^{1/3}$.\index{transcendental number}\index{Liouville's
inequality}\index{algebraic number}

\medskip
\noindent\textbf{Part I --- How well can rationals be
approximated?}
\begin{enumerate}
  \item Show that two distinct rationals $\frac ab \neq \frac pq$
        (written with $b, q \geq 1$) satisfy $\bigl|\frac ab -
        \frac pq\bigr| \geq \frac{1}{bq}$. Deduce: if $x =
        \frac ab$ and $0 < \bigl|x - \frac pq\bigr| <
        \frac{1}{bq}$, no such $\frac pq$ exists --- a rational
        repels all other rationals at scale $\frac 1q$.
  \item Prove that for \emph{every} rational $\frac pq$ ($q \geq
        1$): $\bigl|\sqrt2 - \frac pq\bigr| \geq
        \frac{1}{4q^2}$ \emph{(if the distance exceeds $1$ this
        is clear; otherwise bound $\abs{\sqrt2 + p/q} < 4$ and
        use the nonzero integer $\abs{p^2 - 2q^2} \geq 1$)}.
  \item In the other direction: check that $(p, q) \mapsto (p +
        2q, p + q)$ preserves $\abs{p^2 - 2q^2} = 1$, generate
        from $(1,1)$ the pairs $(3,2)$, $(7,5)$, $(17,12)$,
        $(41,29)$, $(99,70)$, and show each satisfies
        \[
        \Bigl|\sqrt2 - \frac pq\Bigr| =
        \frac{1}{q^2\,(\sqrt2 + p/q)} < \frac{1}{2q^2} :
        \]
        infinitely many approximations of order $2$. With
        question 2: the approximation exponent of $\sqrt 2$ is
        \emph{exactly} $2$.
  \item (Dirichlet) Let $x$ be irrational and $N \in \N^*$.
        Consider the $N + 1$ fractional parts of $0, x, 2x,
        \dots, Nx$ in the $N$ boxes $\intco{\frac kN}{\frac{k +
        1}{N}}$: by the pigeonhole principle
        (\cref{cor:b1:counting:pigeonhole}), two fall in one box.
        Deduce $q \leq N$ and $p$ with $\abs{qx - p} <
        \frac 1N$, hence infinitely many rationals with
        $\bigl|x - \frac pq\bigr| < \frac{1}{q^2}$: \emph{every}
        irrational is approximable to order $2$.
\end{enumerate}

\medskip
\noindent\textbf{Part II --- Liouville's inequality.} Let $x$ be
irrational and algebraic.
\begin{enumerate}[resume]
  \item Show that among the nonzero integer \omterm{def:b1:poly:def}{polynomials}
        vanishing at $x$ there is one, say $P$ of degree $d$,
        with \emph{no rational root}; and check $d \geq 2$
        \emph{(divide out a factor $X - \frac ab$ over $\Q$ and
        clear denominators; degree $1$ would make $x$
        rational)}.
  \item Show that for every rational $\frac pq$ ($q \geq 1$):
        $\bigl|P\bigl(\frac pq\bigr)\bigr| \geq \frac{1}{q^d}$
        \emph{($q^d P(p/q)$ is a nonzero integer)}.
  \item Let $M = \max_{\intcc{x-1}{x+1}} \abs{P'}$
        (\cref{thm:b1:continuity:evt}). Using the mean value
        theorem between $x$ and $\frac pq$, prove
        \emph{Liouville's inequality}: with $C = \min\bigl(1,
        \frac 1M\bigr) > 0$,
        \[
        \Bigl| x - \frac pq \Bigr| \geq \frac{C}{q^{\,d}}
        \qquad\text{for every rational } \frac pq,\ q \geq 1 .
        \]
  \item Call $x$ a \emph{Liouville number} when for every $n \in
        \N$ there is a rational $\frac pq$ with $q \geq 2$ and
        $0 < \bigl|x - \frac pq\bigr| < q^{-n}$. Prove that a
        Liouville number is irrational \emph{(question 1: choose
        $n$ with $2^{\,n-1} > b$)}.
  \item Prove Liouville's theorem: \emph{a Liouville number is
        transcendental} \emph{(combine questions 7 and 8: the
        inequality $C < q^{\,d-n}$ fails for large $n$)}.
\end{enumerate}

\medskip
\noindent\textbf{Part III --- The number $L$.}
\begin{enumerate}[resume]
  \item Let $L$ be the value (in the sense of
        \cref{pb:b1:reals:1}) of the decimal digit string with
        digit $1$ at the positions $n!$ ($n = 1, 2, 3, \dots$)
        and $0$ elsewhere, i.e.\ $L = \sup_k t_k$ with $t_k =
        \sum_{n=1}^{k} 10^{-n!}$. Write out the first $25$
        digits of $L$.
  \item Prove the tail bracketing, for every $k \geq 1$:
        \[
        10^{-(k+1)!} \;\leq\; L - t_k \;\leq\;
        \frac{10}{9}\,10^{-(k+1)!} \;<\; 2\cdot 10^{-(k+1)!}
        \]
        \emph{(bound every partial sum beyond $t_k$ by a finite
        geometric sum)}.
  \item Write $t_k = \frac{p_k}{q_k}$ with $q_k = 10^{k!}$.
        Show $0 < L - \frac{p_k}{q_k} < \frac{2}{q_k^{\,k+1}}$,
        and conclude that $L$ is a Liouville number in the sense
        of question 8.
  \item Conclude: $L$ is transcendental --- the first explicit
        example in history (Liouville, 1844). Cross-check its
        irrationality directly: its digits are not eventually
        periodic (growing gaps, as in \cref{pb:b1:reals:1},
        question 20).
  \item Generalize: replace each digit $1$ by an arbitrary
        nonzero digit $d_n \in \intint{1}{9}$. Show the value is
        still a Liouville number, and deduce --- by the diagonal
        argument of \cref{pb:b1:reals:1} (question 22) applied
        to these digit choices --- that there are uncountably
        many \omterm{pb:b1:logic:1}{transcendental numbers} of this shape.
\end{enumerate}

\medskip
\noindent\textbf{Part IV --- The hierarchy of approximation
orders.} Say $x$ is \emph{approximable to order $\mu$} when for
some constant $c > 0$ infinitely many rationals satisfy
$\bigl|x - \frac pq\bigr| < \frac{c}{q^{\mu}}$.
\begin{enumerate}[resume]
  \item Assemble the hierarchy from Parts I--III: rationals are
        approximable to order $1$ and no better; $\sqrt 2$ to
        order $2$ and no better; every irrational to order at
        least $2$; an \omterm{pb:b1:logic:1}{algebraic number} of degree $d$ to no order
        beyond $d$; Liouville numbers to every order. Justify
        each claim by citing the relevant question.
  \item Show that $L + r$ is a Liouville number for every
        rational $r = \frac ab$ \emph{(translate the
        approximants: the new denominators are $b\,q_k$)}.
        Conclude that Liouville --- hence transcendental ---
        numbers are \omterm{def:b1:topology:dense}{dense} in $\R$.
  \item (Cantor, 1874) Prove that the \omterm{def:b1:logic:sets}{set} of \omterm{pb:b1:logic:1}{algebraic numbers}
        is countable: there are finitely many integer
        \omterm{def:b1:poly:def}{polynomials} with degree plus sum of $\abs{\text{coefficients}}$
        bounded by $h$, each with at most $\deg$ roots; a
        countable union of \omterm{def:b1:counting:card}{finite sets} is countable. Since no
        sequence exhausts $\R$ (\cref{pb:b1:reals:1}, question
        22), \omterm{pb:b1:logic:1}{transcendental numbers} exist --- in fact form an
        uncountable \omterm{def:b1:logic:sets}{set}. Compare the two proofs: what does
        Liouville's give that Cantor's cannot?
  \item Prove directly from question 2 that $\sqrt 2$ is
        \emph{not} a Liouville number \emph{(for $n \geq 3$,
        the inequality $q^{-n} > \frac{1}{4q^2}$ bounds $q$;
        then only finitely many candidate rationals remain, all
        at positive distance from $\sqrt2$)}. Generalize:
        no \omterm{pb:b1:logic:1}{algebraic number} is Liouville.
\end{enumerate}

\medskip
\noindent\textbf{Part V --- Effective constants.}
\begin{enumerate}[resume]
  \item For the Pell pair $(99, 70)$: verify $99^2 - 2\cdot70^2
        = 1$ and evaluate the exact error
        \[
        \sqrt2 - \frac{99}{70}
        = \frac{-1}{70^2\,\bigl(\sqrt2 + \frac{99}{70}\bigr)},
        \qquad
        \Bigl|\sqrt 2 - \frac{99}{70}\Bigr| \approx 7.2\cdot
        10^{-5} :
        \]
        five correct digits from a three-digit fraction.
  \item Run Part II on $x = 2^{1/3}$, $P = X^3 - 2$: check $P$
        has no rational root, bound $M = \max_{\intcc{x-1}{x+1}}
        3t^2 \leq 3\,(1 + 2^{1/3})^2 < 16$, and conclude the
        effective inequality
        \[
        \Bigl| 2^{1/3} - \frac pq \Bigr| \geq
        \frac{1}{16\,q^3} \qquad \text{for all } \frac pq .
        \]
  \item Payoff: show that any rational approximating $2^{1/3}$
        within $10^{-6}$ must have denominator $q \geq 40$.
  \item Show that the base $10$ is irrelevant: the binary
        analogue $\sum_{n\geq1} 2^{-n!}$ (value of the binary
        string with ones at factorial positions) is also a
        Liouville number, hence transcendental.
\end{enumerate}

\medskip
\noindent\textbf{Part VI --- Frontiers and synthesis.}
\begin{enumerate}[resume]
  \item Let $x^\dagger$ be the value of the decimal string with
        ones exactly at the positions $3^k$ ($k \geq 0$). Show
        $x^\dagger$ is approximable to order $3$, and deduce
        from Liouville's inequality that $x^\dagger$ is neither
        rational nor a quadratic irrational. Explain why the
        method stalls there: order $3$ is compatible with
        algebraicity of degree $\geq 3$, and closing that gap
        (any exponent $> 2$ suffices, for every \omterm{pb:b1:logic:1}{algebraic
        number}) is Roth's theorem, far beyond this volume.
  \item Quantify Cantor: show that the \omterm{pb:b1:logic:1}{algebraic numbers} of
        degree $\leq d$ given by \omterm{def:b1:poly:def}{polynomials} with coefficients
        in $\intint{-H}{H}$ number at most $d\,(2H + 1)^{d+1}$.
        (This finiteness is what made question 17 work.)
  \item Synthesis, one sentence each: (i) locate the single
        analytic ingredient of Liouville's proof (which theorem
        of this chapter, used where); (ii) state the tension
        that powers it (integrality forces $\abs{P(p/q)} \geq
        q^{-d}$, smoothness forbids $\abs{P(p/q)} > M\abs{x -
        p/q}$); (iii) contrast Liouville's and Cantor's proofs
        of the existence of \omterm{pb:b1:logic:1}{transcendental numbers}; (iv) name
        where this volume meets the theme again --- the weekend
        problem of \cref{ch:b1:integration} proves $\pi$
        irrational by the same integrality-versus-smallness
        squeeze, with integrals in place of \omterm{def:b1:derivative:def}{derivatives}.
\end{enumerate}
\end{problem}
